% !PRINTED 19
\weberpage{19}

The only number having absolute value $0$ is $0$. Every positive
number occurs as the absolute value of infinitely many complex
numbers, and the image points of all numbers with the same absolute
value lie on a circle whose center is at the origin.

Two imaginary numbers which differ only by the sign of $i$, namely
$x+yi$ and $x-yi$, are called conjugate imaginary. Their product is
the square of the absolute value of either one.

If one of two conjugate imaginary numbers is zero, the other is also
zero; hence in every correct numerical equation one may replace $i$
by $-i$ without disturbing its correctness.

We shall also state the frequently used theorem that the absolute
value of the sum of two nonzero complex numbers is never greater than
the sum of the absolute values of the summands, and is equal to it only
when the ratio (the quotient) of the two summands is real and positive.
Namely, let
\[
 z=x+yi,
 \qquad c=a+bi,
 \qquad Z=(x+a)+(y+b)i,
\]
\[
 \varrho^2=x^2+y^2,
 \qquad r^2=a^2+b^2,
 \qquad R^2=(x+a)^2+(y+b)^2,
\]
then
\[
 (r+\varrho-R)(r+\varrho+R)
   =(r+\varrho)^2-R^2=2(r\varrho-ax-by);
\]
this is certainly positive when $ax+by\leqq 0$. If, however,
$ax+by>0$, then from
\[
 r^2\varrho^2-(ax+by)^2=(ay-bx)^2,
\]
it follows that $r\varrho\geqq ax+by$, with equality only when
$ay-bx=0$. Hence $r+\varrho>R$, except in the special case
$r+\varrho=R$ when $ay-bx=0$ and $ax+by>0$, which proves the
assertion.

In the geometric representation this theorem says that in a triangle
one side is smaller than the sum of the other two. The theorem has the
further consequence that the absolute value of a sum cannot be smaller
than the difference of the absolute values of the summands. For,
applying the preceding theorem to the sum $c=Z-z$, one obtains
$r\leqq R+\varrho$, or
\[
 R\geqq r-\varrho.
\]
%CONT>>
